%!TEX root = ../../csuthesis_research.tex

\chapter{总体方案设计与模块实现}

\section{总体技术框架}

本调研报告所提总体方案以第二章的递归岭回归闭式解为统一核心，把面向四类典型持续学习任务的方法组织成可复用的“三阶段”技术路线。该路线在工程上既保留了深度神经网络对原始信号的表征能力，又用闭式解替换传统反向传播分类头的多轮迭代过程，从而在不引入历史样本回放的前提下，逐阶段地维护可解析更新的判别边界。三阶段的具体内容如下：
\begin{enumerate}
    \item \textbf{初始表征学习阶段}：在第一任务上以常规监督学习方式训练特征编码器，目的是获得迁移性较强、对类别扰动较稳健的中间表征，作为后续阶段的稳定基础~\citep{he2016resnet,chen2017deeplabv3}；
    \item \textbf{解析对齐阶段}：在初始任务训练完成后冻结编码器，用岭回归闭式解构造一次性解析分类头，把表征空间与首任务标签空间显式对齐，避免后续递推累积梯度噪声~\citep{zhuang2022acil,mcdonnell2023ranpac}；
    \item \textbf{增量递推阶段}：对后续到达的任务仅维护逆自相关矩阵 $\Psi_t$ 与互相关矩阵 $\mathbf{Q}_t$，按照第二章式~\eqref{eq:phi_update_report} 完成分类头的递归更新，从而获得与联合训练严格等价的解。
\end{enumerate}

为进一步突出本方案相对于传统基于反向传播的增量学习方法的差异，将两类范式的关键工程属性对照如下：
\begin{description}
    \item[历史样本依赖] 传统 BP 增量方法多依赖经验回放或生成回放以缓解灾难性遗忘~\citep{rebuffi2017icarl,shin2017continual,buzzega2020dark}；本方案仅维护协方差类统计量，不保留任何原始历史样本。
    \item[更新机制] 传统方法在新任务到来时需要数十乃至上百个 epoch 的反向传播迭代；本方案的分类头更新为一次性线性代数运算，不存在迭代轮数随任务数膨胀的问题。
    \item[等价性保证] 正则化方法与梯度约束方法仅在近似意义上保护旧知识~\citep{kirkpatrick2017ewc,li2017learning,lopez2017gradient}；本方案在统计量充分维护的条件下，递推解与联合训练解严格相等。
    \item[超参数敏感性] 传统方法对学习率、批大小、训练轮数等敏感；本方案在编码器冻结后仅剩正则系数 $\gamma$ 一个超参数。
    \item[扩展性] 结构扩展类方法需要为每个任务新增网络分支~\citep{rusu2016progressive,mallya2018piggyback}；本方案在分类头维度上线性扩展，主干结构保持稳定。
\end{description}

从工程实现角度，总体方案可进一步细化为“离线初始化 + 在线递推”两层：
\begin{itemize}
    \item \textbf{离线初始化层}：完成主干网络训练、解析对齐与正则系数搜索，结果以模型权重与初始统计量两类制品形式归档；
    \item \textbf{在线递推层}：对每个新任务仅执行一次特征前向、一次统计量更新与一次解析分类头更新，整个流程不再触发反向传播。
\end{itemize}

\section{统一记号与接口设计}

为使后续 Ch4 中四类任务能够共用同一套数学接口，本节给出统一记号约定，并明确每类任务在该记号体系下对应的具体语义。设第 $t$ 阶段输入特征矩阵为 $\mathbf{X}_t\in\mathbb{R}^{N_t\times d}$，标签矩阵为 $\mathbf{Y}_t$，其中沿用第二章式~\eqref{eq:y_split_report} 的分块形式，把 $\mathbf{Y}_t$ 划分为旧类监督子矩阵 $\bar{\mathbf{Y}}_t$ 与新增类监督子矩阵 $\tilde{\mathbf{Y}}_t$。
\section{模块化设计}

本方案在工程上拟划分为四个相对独立的模块，模块之间以矩阵形式的接口张量相互连接，便于在不同任务上替换或扩展。各模块的设计细节如下。

\subsection{模块一：特征编码模块}

特征编码模块负责把不同模态的原始输入统一压缩为定长向量或位置级特征矩阵。其输入为图像张量、时间序列张量、点云坐标或多模态隐藏状态序列；输出为统一形态的特征矩阵 $\mathbf{X}_t^{\mathrm{encoder}}\in\mathbb{R}^{N_t\times d}$。关键参数包括编码器架构（如 ResNet~\citep{he2016resnet}、一维卷积网络、DGCNN~\citep{wang2019dgcnn}、冻结的多模态大语言模型~\citep{liu2023llava}）以及主干冻结时机。本模块与第二章定理的接口体现在：编码器输出直接作为 $\mathbf{X}_t$ 的来源；初始任务训练完成后该模块进入只读状态，从而保证后续阶段统计量更新的稳定性。本方案拟在第一阶段以监督学习方式完成该模块的训练并冻结，避免后续任务对历史表征产生破坏。

\subsection{模块二：高维映射与增强模块}

该模块设计目标在于补偿冻结编码器引致的可塑性损失。其输入为编码器输出 $\mathbf{X}_t^{\mathrm{encoder}}$；输出为高维特征矩阵 $\mathbf{X}_t^{E}=\mathrm{ReLU}(\mathbf{X}_t^{\mathrm{encoder}}\mathbf{R}^E)$，其中 $\mathbf{R}^E\in\mathbb{R}^{d\times d_E}$ 为随机初始化矩阵。关键参数包括随机种子、扩展维度 $d_E$ 与非线性激活类型。本模块借鉴 RanPAC~\citep{mcdonnell2023ranpac} 与 GKEAL~\citep{GKEAL_Zhuang_CVPR2023} 中关于随机投影提升线性可分性的结论。本模块与第二章定理的接口体现在：其输出直接替换 $\mathbf{X}_t$ 进入递归更新公式，使解析分类头在更高维空间中具备更强的判别能力。本方案拟在解析对齐与递推阶段始终启用该模块，并通过预先固定随机种子保证不同阶段的可复现性。

\subsection{模块三：递归闭式解更新模块}

递归闭式解更新模块是整个方案的算法核心。其输入为当前阶段的高维特征矩阵 $\mathbf{X}_t^{E}$、分块标签矩阵 $(\bar{\mathbf{Y}}_t,\tilde{\mathbf{Y}}_t)$ 以及上一阶段缓存的统计量 $(\Psi_{t-1},\widehat{\Phi}_{t-1})$；输出为更新后的统计量 $\Psi_t$ 与分类头参数 $\widehat{\Phi}_t$。关键参数包括岭回归正则系数 $\gamma$、矩阵求逆所用的数值方法以及统计量的存储精度。本模块严格实现第二章的递推公式~\eqref{eq:psi_update_report} 与~\eqref{eq:phi_update_report}，并复用 Woodbury 恒等式以避免对 $\Psi_{t-1}^{-1}$ 的显式求逆。与第二章定理的接口体现在：本模块即为定理的工程实现，调用方仅需提供四元组 $(\mathbf{X}_t^{E},\bar{\mathbf{Y}}_t,\tilde{\mathbf{Y}}_t,\Psi_{t-1})$ 即可获得新的分类头，无需访问任何历史样本。

\subsection{模块四：任务适配模块}

任务适配模块负责把上述统一接口与具体任务的特殊性进行衔接。其输入为任务相关的辅助信息，如时间序列的多尺度特征、分割任务的伪标签判定阈值、多模态路由的专家编号映射等；输出为符合统一接口形态的特征与标签张量。关键参数包括：医疗任务的类别增量切分协议、时序任务的多尺度层索引集合、分割任务的伪标签置信度阈值 $\tau$、多模态路由任务的池化粒度。该模块与第二章定理的接口体现在：在不修改递推公式的前提下，仅通过调整 $\mathbf{X}_t$ 与 $\mathbf{Y}_t$ 的构造方式即可适应不同任务，因而保证算法核心在四类应用上保持一致。

\section{工程流程}

为保证本调研报告所提方法可复现，给出统一的六步工程流程，并标注每步的输入、输出与计算复杂度（设 $N_t$ 为阶段样本数，$d$ 为编码器输出维度，$d_E$ 为随机高维映射维度）：
\begin{enumerate}
    \item \textbf{数据读取与任务划分}。\textit{输入}：原始数据集与任务切分协议；\textit{输出}：阶段化数据流 $\{\mathcal{D}_t\}_{t=1}^{T}$；\textit{复杂度}：与数据集规模线性相关。
    \item \textbf{初始任务训练与主干冻结}。\textit{输入}：$\mathcal{D}_1$；\textit{输出}：冻结的编码器权重；\textit{复杂度}：与常规深度学习训练相当，是整个流程中唯一显著的反向传播开销。
    \item \textbf{特征抽取与随机高维映射}。\textit{输入}：阶段样本 $\mathbf{X}_t$；\textit{输出}：高维特征 $\mathbf{X}_t^{E}$；\textit{复杂度}：$\mathcal{O}(N_t d d_E)$，每个样本仅做一次前向。
    \item \textbf{递归统计量更新与闭式解分类头更新}。\textit{输入}：$(\mathbf{X}_t^{E},\bar{\mathbf{Y}}_t,\tilde{\mathbf{Y}}_t,\Psi_{t-1},\widehat{\Phi}_{t-1})$；\textit{输出}：$(\Psi_t,\widehat{\Phi}_t)$；\textit{复杂度}：矩阵乘法主导，$\mathcal{O}(d_E^2 N_t + d_E^3)$，与任务数 $t$ 无关。
    \item \textbf{阶段评测与日志归档}。\textit{输入}：累计测试集与当前 $\widehat{\Phi}_t$；\textit{输出}：阶段精度、遗忘度量与日志；\textit{复杂度}：仅一次前向推理，开销极小。
    \item \textbf{进入下一任务循环}。仅保留 $(\Psi_t,\widehat{\Phi}_t)$ 与冻结编码器进入下一轮，历史样本与中间梯度一律丢弃。
\end{enumerate}

上述流程在工程上具备良好的可拆分性：第 2 步可借助 GPU 集群完成离线训练，第 3 至第 5 步在新任务到来时可以以服务化方式即时执行。该设计与持续学习对系统资源弹性扩展的要求~\citep{wang2024survey} 相吻合。

\section{方案与持续学习关键约束的对应关系}

调研报告第一章梳理了持续学习在工程落地中面临的四类关键约束，本节逐一说明本方案的应对机制。
\begin{itemize}
    \item \textbf{数据约束}。许多实际场景因隐私、版权或存储原因不允许保留历史样本~\citep{de2021continual,wang2024survey}。本方案在递推阶段仅维护协方差类统计量 $\Psi_t$ 与 $\mathbf{Q}_t$，不存储任何历史原始数据，天然契合无回放需求。
    \item \textbf{算力约束}。增量阶段的算力预算通常远低于初始训练阶段，特别是在端侧或医院本地服务器上更为突出~\citep{souza2020challenges}。本方案在该阶段仅执行一次矩阵运算，不再进行反向传播，因而对算力预算友好。
    \item \textbf{稳定性约束}。稳定性—可塑性两难是持续学习的核心矛盾~\citep{kirkpatrick2017ewc,li2017learning}。本方案通过冻结主干并以统计量形式显式记忆历史信息，使得分类边界仅由累计统计量决定，从而在理论上与联合训练等价，确保稳定性。
    \item \textbf{扩展性约束}。新任务的引入不应导致整体模型规模或推理延迟显著膨胀~\citep{rusu2016progressive,mallya2018piggyback}。本方案仅在分类头维度上线性扩展，主干和高维映射模块保持不变，因而具备良好的长期扩展能力。
\end{itemize}

\section{本方案与现有方法的差异}

现有持续学习方法大致可归为回放方法、正则化方法与结构扩展方法三类~\citep{wang2024survey,de2021continual}。回放方法通过经验回放或生成回放重现历史样本~\citep{rebuffi2017icarl,shin2017continual,buzzega2020dark}；正则化方法以惩罚项约束参数偏移~\citep{kirkpatrick2017ewc,li2017learning,aljundi2018memory}；结构扩展方法为每个任务新增子网络或专家~\citep{rusu2016progressive,mallya2018piggyback,wang2022dualprompt}。本方案在以下几方面与上述三类方法形成显著差异：
\begin{itemize}
    \item 与回放方法相比，本方案完全不保留历史样本，仅维护统计量，避免了隐私泄露与缓冲区大小受限的问题；
    \item 与正则化方法相比，本方案在统计量充分维护时与联合训练严格等价，不存在因近似而引入的遗忘上界；
    \item 与结构扩展方法相比，本方案仅在分类头维度上做线性扩展，主干保持冻结，避免了任务数增多时模型规模的非线性膨胀；
    \item 本方案在四类任务上提出统一的接口设计，并明确给出医学、时序、分割与多模态路由四类场景的适配机制。
\end{itemize}

\section{本章小结}

本章在第二章统一数学框架基础上，给出了面向四类持续学习任务的总体方案设计与模块化实现。本方案以三阶段技术路线为骨架，以四模块结构为支撑，以六步工程流程为操作规范，并通过统一记号与接口设计实现了不同任务之间的算法核心复用。本方案对数据、算力、稳定性与扩展性四类关键约束逐一给出了应对机制，并与回放、正则化、结构扩展三类主流方法形成了清晰的差异。第四章拟在本设计基础上，分别在医学影像疾病分类、时间序列分类、语义分割与多模态专家路由四类典型任务上完成具体方法的实例化，并通过实验进一步验证本方案的有效性与可扩展性。
